\chapter{Common emitter amplifier}
This experiment explores the usage of a transistor as an amplifier.

\section{Single stage}
\begin{figure}[tbp]
	\centering
	%\includegraphics[width=.5\textwidth]{./img/com_emitter_amp.pdf}
	\caption[Single Stage Common Emitter Amplifier]{Single stage common emitter amplifier}
	\label{fig:com_emitter_amp}
\end{figure}
The circuit shown in \autoref{fig:com_emitter_amp} acts as a voltage amplifier.
The input and output capacitors remove any constant component of the input.
\comp{R_1} and \comp{R_2} are used to bias the transistor, so that it remains in active mode during operation.
This method of biasing vastly reduces the effects of varying $\beta$ as a cause of temperature drift by holding the base bias at a constant voltage.
\comp{R_E} is used to compensate for temperature effects by implementing a negative feedback, albeit sacrificing a small amount of amplification, depending on the conductance $S=-\frac{I_\text{C,bias}}{U_\text{T}}$, where $U_\text{T}$ denotes the temperature voltage ($\approx\SI{26}{\milli\volt}$).
The temperature dependency of the bias point results out of the fact that the base-emitter voltage \comp{V_BE} decreases by $\SI{2}{\milli\volt\per\kelvin}$ at constant collector current \comp{I_C}.
If the input voltage is kept constant, this temperature increase acts as a raise of the input voltage.
Consequently, temperature drift can be calculated as

\begin{equation}
	\frac{\d U_\text{a}}{\d T}=\frac{\d U_\text{a}}{\d U_\text{e}}\frac{\d U_\text{a}}{\d T}\approx A\cdot\SI{2}{\milli\volt\per\kelvin},
\end{equation}

where $U_\text{a}$ denotes the output voltage, $U_\text{e}$ the input voltage and $T$ the temperature.
This means that temperature drift is proportional to the voltage amplification $A$ in a common emitter amplifier.
Therefore, temperature drift can be lowered by decreasing the amplification.
As temperature drift happens slowly, DC amplification can be decreased, whereas AC amplification can remain unchanged.
This can be realized by putting a capacitor \comp{C_E} in parallel to the emitter resistor \comp{R_E} which will decrease the negative feedback with growing frequency.

The quiescent base voltage is determined by the voltage divider formed by \comp{R_1} and \comp{R_2}:
\begin{equation*}
	V_\text{bias}=V_\text{cc}\cdot\frac{\comp{R_2}}{\comp{R_1}+\comp{R_2}}.
\end{equation*}
